
For convenience, we record the coefficient-by-coefficient interpretation of Theorem~\ref{t:refinedcount}. Let $\Semi_{\bfc,\bfe}(\sigma)$ denote the set of semistandard set-valued fillings of $\sigma$ of content $\bfc$ and excess $\bfe$.

\begin{theorem} \label{t:refinedcount'} Let $\sigma$ be any connected skew shape, and fix sequences $\bfc$ and $\bfe$. Then 
\[
|\Semi_{\bfc,\bfe}(\sigma)| = \sum_{\mu \supseteq \sigma} (-1)^{|A(\mu/\sigma)|} \cdot a_{\sigma,\mu,\bfe} \cdot |\Semi_{\bfc, \mathbf{0}} (\mu)|,
\]
where $a_{\sigma,\mu,\bfe}$ are the nonnegative integers defined in Theorem~\ref{t:refinedcount}.
\end{theorem}

Thus Theorems ~\ref{t:refinedcount} and~\ref{t:refinedcount'} are equivalent.

\begin{rema}
The change from $B(\mu/\sigma)$ in Theorem~\ref{t:refinedcount} to $A(\mu/\sigma)$ in Theorem~\ref{t:refinedcount'} is not accidental; it arises from the definition of $RG$ as a \emph{signed} generating function for set-valued semistandard tableaux.
\end{rema}

Then, by specializing to $\bfw = 1$ in Theorem~\ref{t:refinedcount}, we obtain Theorem~\ref{thm:comb}.



\begin{rema} \label{rem:connection-to-lenart}
Consider the row-bounded, reverse row-strict tableaux of shape $B(\mu/\sigma)$, as in~\eqref{defi3.2_2} above. There is a bijection between this set and the set of row-bounded, row- and column- strictly-decreasing tableaux of shape $B(\mu/\sigma)$, obtained by replacing label $T(i,j)$ with $i-T(i,j)$. Therefore, when $\sigma$ is a straight shape whose highest row is in row $1$, Theorem~\ref{thm:comb} reduces to~\cite[Theorem~2.2]{lenart-combinatorial}. In particular, $A(\mu/\sigma)$ is always empty in this case. 

We also note that when $N=|\sigma|+1$, Theorem~\ref{thm:comb} specialized to the monomial $x_1\cdots x_N$ is equivalent to~\cite[Theorem~2.8]{clpt}.
Moreover a proof using a row insertion algorithm in the special case that $\sigma$ is a straight shape and $N=|\sigma|+1$ is presented in~\cite[Proposition~3.9]{reiner-tenner-yong}.
\end{rema}




\begin{rema}\label{rem:connection-to-act}
The determinantal formula of~\cite{anderson-chen-tarasca-k-classes} can also be expanded as a similar sum involving enumeration of standard Young tableaux on larger skew shapes (see~\cite[Theorem~C]{anderson-chen-tarasca-k-classes}).
Thus, Theorem~\ref{thm:comb} establishes in a purely combinatorial manner that the determinantal formula in~\cite{anderson-chen-tarasca-k-classes} is equal to the number of set-valued tableaux.
\end{rema}

\begin{rema}\label{rem:darij}
D.~Grinberg has pointed out the row-refined skew stable Grothendieck polynomials $RG_\sigma(\mathbf{x};\mathbf{w})$ in Definition~\ref{def:rg}, restricted to straight shapes $\sigma$, are Hall-dual to the power series $\tilde{g}_{\mu}(\mathbf{x};\mathbf{w})$ defined in~\cite{galashin-grinberg-liu-refined}. In other words, 
\[
\langle RG_\sigma(\mathbf{x};\mathbf{w}), \tilde{g}_\mu(\mathbf{x};\mathbf{w})\rangle = \delta_{\sigma,\mu}
\]
for straight shapes $\sigma$ and $\mu$. Here $\langle\cdot,\cdot\rangle$ denotes the Hall inner product on symmetric functions in variables $\mathbf{x}$, treating the $w_i$ as scalars. 
\end{rema}






Now we prove Theorem~\ref{t:refinedcount'} using a new generalized row insertion algorithm. This proof occupies the rest of the section. This algorithm extends the set-valued insertion algorithm in~\cite{bandlow-morse-combinatorial} to the case of skew shapes. 

\begin{defi}[RSK row insertion]\label{def:rsk} First, recall the row insertion operation, the atomic operation of the RSK algorithm~\cite[\S~7.11]{ec2} (we present a very slightly more general version). Suppose $\sigma$ is a skew or straight shape and $T$ is a semistandard tableau of shape $\sigma$. Given $k\in \NN$ and $i$, the operation $T\leftarrow_i k$ inserts $k$ in the leftmost box of row $i$ labeled $j>k$, or a new box at the right end of the $i^\text{th}$ row if no box in that row is labeled $>k$ (in the case where there are not yet any boxes in that row, the new box is placed directly below the leftmost entry in the previous row). In the latter case the operation terminates. In the former, we insert $j$ into the $(i+1)^\text{st}$ row of $\sigma$ in the same manner, and repeat down the rows of $\sigma$. The \emph{insertion path} is the sequence of boxes $b_{i,j_1},b_{i+1,j_2},\ldots$ in which insertions occurred; one can check that $j_1 \ge j_2 \ge \cdots$~\cite[Lemma~7.11.2]{ec2}.
\end{defi}

In particular, row insertion inputs a semistandard tableau of shape $\sigma$ and outputs a semistandard tableau of shape $\sigma'$ obtained by adding one box to $\sigma$.

\begin{rema}
Notice that row insertion may be applied without changes to set-valued tableaux in the following situation. Suppose $T$ is a set-valued semistandard tableau of shape $\sigma$. Suppose $k$ is a label in a box $b$ with at least one other label; let $i$ index the row containing $b$. Suppose further that every box in row $i+1,i+2,\ldots$ is labeled with a singleton set. Then one may define the operation $T\leftarrow_i k$ as before, deleting $k$ from box $b$ and row-inserting it in the next row, and repeating. Simply put, the row insertion path does not traverse any box with more than one label in this case.

This observation allows for the next algorithm.
\end{rema}

\begin{enonce}{Algorithm}\label{algorithm}
The \emph{skew set-valued} row insertion algorithm, for a connected skew shape $\sigma$, is as follows. The input is

\begin{enumerate}
\item\label{algo3.12_1_1} a skew shape $\lambda \supseteq \sigma$ with $B(\lambda/\sigma)=\emptyset$,
\item\label{algo3.12_1_2} $T'$ a reverse-row-strict, row-weakly-bounded tableau on $A(\lambda/\sigma)$, and
\item\label{algo3.12_1_3} $T\in \Semi_{\bfc,\bfe-c(T')}(\lambda)$.
\end{enumerate}

The output will be:
\begin{enumerate}
\item\label{algo3.12_2_1} a skew shape $\mu\supseteq \sigma$ with $A(\mu/\sigma)=\emptyset$,
\item\label{algo3.12_2_2} $T''$ a reverse-row-strict, row-bounded tableau on $B(\mu/\sigma)$ with $c(T'')=\bfe$, and
\item\label{algo3.12_2_3} $\wt{T} \in \Semi_{\bfc,\bfzr}(\mu)$.
\end{enumerate}
\end{enonce}

The algorithm proceeds as follows. Let $r$ be the number of rows of $\sigma$. For each $k=r,\ldots,1$ (in descending order), we will do two ``sweeps'' of $\sigma$. First we sweep out all labels in row $k$ that are not the minimum in their box, via row-inserting them downward. Then we sweep out all labels in all (singly-labeled) boxes $b$ for which $T'(b) = k$, again via row insertion. These boxes need not be in row $k$. In the auxiliary labeling $T''$, the newly created boxes are labeled $k$, and properties of row insertion will imply that at most one box in each column of $T''$ is labeled $k$. An example is given in Example~\ref{ex:RSK}. 

Now we describe the algorithm more precisely. For $k=r,\ldots,1$, proceed as follows. First, let $m$ be the maximum label in the rightmost box of row $k$ that has multiple labels. Delete $m$ and insert $m$ into the leftmost box of row $k+1$ labeled $m_2>m$, or a new box at the right end of the $(k+1)^\text{st}$ row if no box in that row is labeled $>m$. In the latter case the operation terminates; the new box is labeled $k$ in the auxiliary filling $T''$. In the former, we insert $m_2$ into the $(k+2)^\text{nd}$ row of $\sigma$ in the same manner, and repeat down the rows of $\sigma$. This is the familiar row-insertion operation of Definition~\ref{def:rsk}. The \emph{insertion path} is the sequence of boxes $b_0 = (k,j_0),b_1=(k\!+\!1,j_1),\ldots$ in which insertions occurred; one can check that $j_0 \ge j_1 \ge \cdots$ (\cite[Lemma~7.11.2]{ec2}). Repeat row-insertion on the maximum label in the rightmost non-singly valued box in row $k$, until that row has only singly-valued boxes. 

The second part of step $k$ is as follows. Since $T'$ is row-weakly-bounded and reverse row-strict, it follows that there is at most one box $(i,j)\in A(\lambda/\sigma)$ in each row such that $T'(i,j)= k$; furthermore $i\ge k$ if so, so that $T(i,j)$ must consist of a single label. So for each such box $(i,j)$, taken in order with $i$ increasing, delete the box and row-insert the label $T(i,j)$ it starting in row $i+1$. When the operation terminates, the new box is labeled $k$ in the auxiliary filling $T''$. We note for later use that in this stage, every box labeled $k$ in $T''$ is the leftmost of its row, since all boxes labelled $>k$ have already been removed.



\begin{exam}\label{ex:RSK}
Let 
\[
%\abovedisplayskip0pt
\belowdisplayskip0pt
\sigma = \text{\tiny \young(::\ ,::\ ,:\ \ ,\ \ )} \qquad \lambda = \text{\tiny \young(:\ \ ,\ \ \ ,\ \ \ ,\ \ )} \qquad T' = \text{\tiny \young(:1,21,1)} \qquad T=
\text{\tiny\ytableausetup{boxsize=.75cm} 
\begin{ytableau}
 \none & 2 &3\\
 1,4 & 6 & 8 \\
 5,7 & 9 & 10,\blue{13} \\
 11 & 12
 \end{ytableau}}.
\]
The algorithm gives
\[
\begin{split}
\ytableausetup{boxsize=.48cm} 
\text{\tiny \begin{ytableau}
 \none & 2 &3\\
 1,4 & 6 & 8 \\
 5,\blue{7} & 9 & 10 \\
 11 & 12 & \red{13}
 \end{ytableau}}
 \qquad
\text{\tiny\begin{ytableau}
 \none & 2 &3\\
 1,\blue{4} & 6 & 8 \\
 5 & 9 & 10 \\
 7 & 12 & \red{13}\\
 \red{11}
 \end{ytableau}}
 \qquad
\text{\tiny\begin{ytableau}
 \none & 2 &3\\
 \blue{1} & 6 & 8 \\
 4 & 9 & 10 \\
 5 & 12 & \red{13} \\
 \red{7}\\
 \red{11}
 \end{ytableau}}
 \qquad
\text{\tiny\begin{ytableau}
 \none & \blue{2} &3\\
 \none & 6 & 8 \\
 1 & 9 & 10 \\
 4 & 12 & \red{13} \\
 \red{5}\\
 \red{7}\\
 \red{11}
 \end{ytableau}}
\\
\ytableausetup{boxsize=.48cm} 
\text{\tiny \begin{ytableau}
 \none & \none &3\\
 \none & \blue{2} & 8 \\
 1 & 6 & 10 \\
 4 & 9 & \red{13} \\
 \red{5} & \red{12}\\
 \red{7}\\
 \red{11}
 \end{ytableau}}
 \qquad
\text{\tiny \begin{ytableau}
 \none & \none &3\\
 \none & \none & 8 \\
 \blue{1} & 2 & 10 \\
 4 & 6 & \red{13} \\
 \red{5} & \red{9}\\
 \red{7}& \red{12}\\
 \red{11}
 \end{ytableau}}
\qquad
\text{\tiny \begin{ytableau}
 \none & \none &3\\
 \none & \none & 8 \\
 \none & 2 & 10 \\
 1 & 6 & \red{13} \\
 \red{4}& \red{9}\\
 \red{5} & \red{12}\\
 \red{7}\\
 \red{11}
 \end{ytableau}}
\end{split}
 \]
and the auxiliary tableau $T''$ is 
\[
\young(::3,31,21,2,1)
\]
\end{exam}


\begin{lemma}\label{lem:asclaimed}
The output of Algorithm~\ref{algorithm} takes the claimed form.
\end{lemma}

\begin{proof}
We remark that the process in Algorithm~\ref{algorithm} preserves the property that \emph{every box in row $k+1$ and below has exactly one label in it}, so the row-insertion is always well-defined. The process also clearly preserves the content of the tableau $T$. Thus iterating the described two-step process for $k=r,\ldots,1$ produces the output data $\mu, T'',$ and $\wt T$, with $c(\wt{T}) = c(T)$. Furthermore $T''$ is row-bounded since $T'$ was row-weakly-bounded. To conclude that the output is as claimed, the only thing left to show is that the labeling $T''$ of $B(\mu/\sigma)$ is reverse row-strict.

 Indeed, since the rows are processed in the order $r,\ldots,1$ in Algorithm~\ref{algorithm}, it is enough to show that for a fixed $k\in\{1,\ldots,r\}$ that no two boxes labelled $k$ in $T''$ lie in the same row. This follows from the standard fact that row-insertion paths move weakly to the left. Precisely: Suppose $m$ and $m'$ are labels that are processed consecutively in step $k$. Let $b_0,b_1,\ldots b_M$ be the insertion path of $m$. By assumption, after $m$ is inserted, every box $b_i$ except possibly $b_0$ is still singly labeled, and $\max(T(b_0)) < T(b_1) < \cdots < T(b_M)$. 

Furthermore, we claim that the label $m'$ is on or to the left of the insertion path of $m$. Indeed, if $m'$ is also in row $k$, then this is clear since $m'<m$; otherwise, we simply note that $m'$ is in the leftmost box of its row, so the claim is also clear. Finally, row-insertion of $m'$ preserves the property of being weakly left of the insertion path of $m$. So the insertion path of $m'$ cannot end to the right of that of $m'$; thus it ends below that of $m'$. This concludes the proof of the lemma.
\end{proof}

Now we show that all possible outputs are attained bijectively by the algorithm.

\begin{prop} \label{p:main}
For any connected skew shape $\sigma$ and any $\bfc$ and $\bfe$, Algorithm~\ref{algorithm} produces a bijection between choices of 
\begin{enumerate}
\item a skew shape $\lambda \supseteq \sigma$ with $B(\lambda/\sigma)=\emptyset$,
\item $T'$ a reverse-row-strict, row-weakly-bounded tableau on $A(\lambda/\sigma)$, and
\item $T\in \Semi_{\bfc,\bfe-c(T')}(\lambda)$;
\end{enumerate}

\noindent
and choices of
\begin{enumerate}
\item a skew shape $\mu\supseteq \sigma$ with $A(\mu/\sigma)=\emptyset$,
\item $T''$ a reverse-row-strict, row-bounded tableau on $B(\mu/\sigma)$ with $c(T'')=\bfe$, and
\item $\wt{T} \in \Semi_{\bfc,\bfzr}(\mu)$.
\end{enumerate}
Therefore, 
\begin{equation}\label{eq:ss}
\sum_{(\lambda, T')} |\Semi_{\bfc, \bfe-c(T')}(\lambda)| = \sum_{(\mu,T'')} |\Semi_{\bfc, \mathbf{0}}(\mu)|
\end{equation}
where
\begin{itemize}
\item the left hand sum ranges over all $\lambda\supseteq \sigma$ with $B(\lambda/\sigma)=\emptyset$, together with a reverse row-strict, row-weakly-bounded labeling $T'$ of $A(\lambda/\sigma)$, and
\item the right hand sum ranges over all $\mu\supseteq \sigma$ with $A(\mu/\sigma) = \emptyset$, together with a reverse row-strict, row-bounded labeling $T''$ of $B(\mu/\sigma)$.
\end{itemize}
\end{prop}

\begin{proof}
\looseness-1
 The skew set-valued row-insertion algorithm in Algorithm~\ref{algorithm} constructs a map
\begin{equation}\label{eq:ssmap}
\coprod_{(\lambda, T')} \Semi_{\bfc, \bfe-c(T')}(\lambda) \xrightarrow{\text{RSK}_\sigma}\coprod_{(\mu,T'')} \Semi_{\bfc, \mathbf{0}}(\mu),
\end{equation}
where the conditions on $\lambda,\mu,T',$ and $T''$ are as in the statement of the proposition. We claim this map is a bijection, and it suffices to provide an inverse. The inverse may be described algorithmically as follows. Given $\mu, T'',$ and $\wt T$ satisfying conditions~\eqref{algo3.12_2_1}, \eqref{algo3.12_2_2}, and~\eqref{algo3.12_2_3} described as the output of Algorithm~\ref{algorithm}, perform the following procedure for $k=1,\ldots, r$. Consider the boxes of $B(\mu/\sigma)$ labelled $k$ in $T''$, in order from highest to lowest row number (\ie lowest to highest on the page). For each such box $b$, delete $b$ and inverse-row-insert its label $m$ upwards, stopping if it reaches row $k$. If the label $m$ lands in a new box $b'$, necessarily in row $\ge k$, then set $T'(b) = k$. An example is given in Example~\ref{ex:RSK}, read in reverse.

The resulting tableau $T'$ is reverse-row-strict by an argument analogous to Lemma~\ref{lem:asclaimed}. So the result of this procedure is the data $\lambda,T',$ and $T$ satisfying the conditions~\eqref{algo3.12_1_1}, \eqref{algo3.12_1_2}, and~\eqref{algo3.12_1_3} described as the input of Algorithm~\ref{algorithm}. 
Now it is evident that the procedure described is in fact inverse to the RSK map in Algorithm~\ref{algorithm}, since each upwards insertion operation is inverse to row insertion, and it processes boxes in the reverse order.
\end{proof}

Now we state a definition and lemma, which will be used to prove Theorem~\ref{t:refinedcount}. We will postpone its proof until after the the proof of Theorem~\ref{t:refinedcount}. The need for Definition~\ref{def:acyclic} as a hypothesis for Lemma~\ref{lem:Gseq} was pointed out to us by D.~Grinberg.

\begin{defi}\label{def:acyclic}
Let $P$ be any finite poset, with its set of cover relations $C = \{(x,y)\in P\times P: x\lessdot y\}$ partitioned into two disjoint sets $C=G\sqcup B$ (called \emph{good} and \emph{bad}, colloquially). We will say that the partition is \emph{acyclic} if the directed graph on the Hasse diagram of $P$ obtained by orienting all good cover relations up and all bad cover relations down is acyclic.
\end{defi}

\begin{exam}
The partition $G\sqcup B$ of the cover relations in the Hasse diagram is drawn below is \emph{not} acyclic.
\[
\xymatrix@R=6mm@C=6mm{&\bullet \ar@{-}[dl]_G\ar@{-}[dr]^B&\\ \bullet \ar@{-}[dr]_G && \bullet \ar@{-}[dl]^B\\ &\bullet&}
\]
\end{exam}


\begin{lemma}\label{lem:Gseq}
Let $P$ be any finite poset, 
with its set of cover relations $C = \{(x,y)\in P\times P: x\lessdot y\}$. Let $C=G\sqcup B$ be an acyclic partition of $C$, with elements of $G$ and $B$ called \emph{good} and \emph{bad} throughout. 
Say that an increasing sequence
\[
\mathcal{I} = (\emptyset = I_0 \subsetneq \cdots \subsetneq I_\ell = P)
\]
of order ideals $I_i$ is a \emph{$G$-sequence} if for every $i=1,\ldots,\ell$, the only cover relations within $I_i \setminus I_{i-1}$ are in $G$. Precisely: if $x,y\in I_i \setminus I_{i-1}$ and $x\lessdot y$ then $(x,y)\in G$. The length of such a $G$-sequence $\mathcal{I}$ is defined to be $|\mathcal{I}|=\ell$.
Then 
\begin{equation}\label{eq:signedcount}
\sum_\text{$\mathcal{I}$ a $G$-sequence} \! (-1)^{|\mathcal{I}|} = \begin{cases}
(-1)^{|P|} & \text{if }G=\emptyset,\\
0 &\text{otherwise.}
\end{cases}
\end{equation}
\end{lemma}

\begin{exam}
Suppose $P = P(\lambda)$ is the poset of boxes of a diagram $\lambda$. If $G=\emptyset$ then the $G$-sequences are in natural correspondence with increasing tableaux of shape $\lambda$ with label set $\{1,\ldots,N\}$ for some $N$. Lemma~\ref{lem:Gseq} states that counting these increasing tableaux, with sign according to the parity of $N$, is $(-1)^{|P|}$. For example, if 
%$P=P(\tiny\Yvcentermath1\yng(3,1))$ 
$P=P({\ytableausetup{boxsize=2.3mm,centertableaux}\ytableaushort{&\\
\none}})$
then the lemma states that 
%\[
%(-1)^4 = \#\left\{\Yvcentermath1\young(123,4),\young(124,3), \young(134,2) \right\} - \#\left\%{\Yvcentermath1\young(123,3) ,\young(123,2)\right\}.
%\]
\[
(-1)^4 = \#\left\{{\ytableausetup{boxsize=3.8mm,centertableaux}\ytableaushort{123,4},\ytableaushort{124,3}},\ytableaushort{134,2} \right\} 
- \#\left\{{\ytableausetup{boxsize=3.8mm,centertableaux}
\ytableaushort{123,3} ,\ytableaushort{123,2}
}\right\}.\]
\end{exam}

Postponing the proof of Lemma~\ref{lem:Gseq}, we now prove Theorem~\ref{t:refinedcount'}. 

